% III. Особености на програмната реализация
% Ориентировъчен обем: 12–13 стр.
%
% Тук влиза и експерименталната част (§ проверка и § резултати): проверката на
% реализацията е нейна особеност, а не отделна глава. Методиката стои ПРЕДИ
% резултатите, защото числото „нула нарушения“ не значи нищо, докато не е казано
% какво се брои за нарушение и как е измерено.
\section{Особености на програмната реализация}
\label{sec:implementation}

\subsection{Организация на хранилището}
\label{subsec:repo}

Хранилището съдържа две изпълними части в обхвата на тази работа, плюс един
отхвърлен прототип, запазен за справка \tabref{tab:repo}. Обемите са дадени в редове
изходен текст и служат за преценка на съотношението между частите, а не като мярка за
вложен труд. Хранилището съдържа и отделно мобилно приложение (\code{mobile/}), извън
обхвата на настоящата работа (Увод, §\ref{sec:introduction}).

\begin{table}[htbp]
\centering
\caption{Съставни части на хранилището в обхвата на работата}
\label{tab:repo}
{\fontsize{10}{12}\selectfont
\begin{tabular}{@{}L{2.3cm}L{7.6cm}r@{}}
\toprule
\textbf{Директория} & \textbf{Съдържание} & \textbf{Редове} \\
\midrule
\code{server/} & услуга на Node.js: приложен програмен интерфейс, връзка с базата, изпращане на поща, изпитания & 929 \\
\code{website/} & уеб клиент без стъпка на компилация (JavaScript) & 4429 \\
\code{website/} & стилове (CSS) & 2029 \\
\code{website/} & разметка на осемте страници (HTML) & 837 \\
\code{web/} & отхвърлен прототип на React с Vite (§\ref{subsec:webtech}), запазен за справка, не се обслужва & 1232 \\
\bottomrule
\end{tabular}}
\end{table}

Наборът с товарните таблици се съхранява на две места в един и същ вид: като
\code{server/data/loads.json} за нуждите на приложния програмен интерфейс и като
\code{website/data.js} с обем 407{,}7 kB за нуждите на уеб клиента, където е обвит в
едно присвояване на глобална променлива, за да бъде зареден с обикновен елемент
\code{script} без система от модули.

\subsection{Реализация на услугата}
\label{subsec:server}

Услугата е изградена като последователност от междинни обработчици, чийто \emph{ред}
носи проектния смисъл. Листинг~\ref{lst:mount} показва решаващата част: публичните
маршрути се монтират преди проверката за връзка с базата от данни, а всички
останали след нея.

\begin{lstlisting}[caption={Ред на монтиране: публичните маршрути преживяват отказ на базата},label={lst:mount}]
app.get("/api/health", (req, res) => res.json({ ok: true, db: isConnected() }));

// Public: load tables - no DB needed, so it's mounted before the DB guard
// (works even if Atlas is down).
app.use("/api/loads", require("./routes/loads"));

// Degrade gracefully: if the DB never connected, return 503 (not a crash).
app.use("/api", (req, res, next) => {
  if (!isConnected()) return res.status(503).json({ error: "Database unavailable" });
  next();
});

app.use("/api/auth", authLimiter, require("./routes/auth"));
app.use("/api/projects", require("./routes/projects"));
\end{lstlisting}

Същото решение е прокарано и в реда на стартиране: услугата започва да слуша по
мрежата \emph{преди} да е установена връзка с базата от данни, а не след това. Ако
редът беше обратният, бавна или неуспешна връзка към отдалечената база би направила
целия сайт недостъпен, включително калкулатора, който не се нуждае от нея.

Две допълнителни решения в услугата заслужават отбелязване, защото не следват от
шаблон, а от наблюдавано поведение.

\textbf{Записване на всяка заявка без тялото ѝ.} Дневникът съдържа метод, път, код на
отговора, продължителност и съобщението за грешка при неуспех, но никога тялото на
заявката или на успешния отговор, тъй като те съдържат пароли и маркери.

\textbf{Възстановяване при отказ на разрешаването на имена.} Свързването към
отдалечената база използва адрес от вида \code{mongodb+srv://}, който изисква справка
за услуга в системата за имена. Много мрежи на доставчици и корпоративни мрежи
отказват тази справка. При разпознаване на такъв отказ услугата повтаря свързването
през публичен разрешител, вместо да приключи с грешка, чието съобщение сочи в грешна
посока (към списъка с разрешени адреси на базата, а не към разрешаването на имена).

\subsection{Удостоверяване и защита}
\label{subsec:security}

Паролите се съхраняват само като хеш, получен с \code{bcrypt} с десет кръга,
което е адаптивна функция, проектирана точно за тази цел~\cite{provos1999bcrypt}.
Проверката при вход използва сравнение на хешове, а отговорът при неуспех е един и
същ независимо от това дали профилът съществува, за да не позволява установяване на
съществуващи адреси.

Сесията е подписан маркер по RFC 7519~\cite{rfc7519} със срок седем дни. Двойният
транспорт, обоснован в §\ref{subsec:authchoice}, е реализиран на едно място
(Листинг~\ref{lst:auth}). Функцията предпочита заглавието пред бисквитката, което
позволява едно и също изпитание (§\ref{subsec:method}) да провери и двата пътя, без
да поддържа две отделни реализации на проверката.

\begin{lstlisting}[caption={Приемане на маркера от заглавие или от бисквитка},label={lst:auth}]
function tokenFrom(req) {
  const auth = req.headers.authorization || "";
  if (auth.startsWith("Bearer ")) return auth.slice(7).trim();
  return (req.cookies && req.cookies[COOKIE]) || null;
}

module.exports = function requireAuth(req, res, next) {
  const uid = verify(tokenFrom(req));
  if (!uid) return res.status(401).json({ error: "Not authenticated" });
  req.userId = uid;
  next();
};
\end{lstlisting}

Бисквитката се поставя с признаци \code{httpOnly} (недостъпна от скрипт на
страницата), \code{sameSite=lax} (не се изпраща при заявки, породени от чужд сайт,
което покрива основния случай на подправяне на заявка между сайтове) и \code{secure}
в производствена среда~\cite{rfc6265}. Тайната за подписване се чете от променлива на
средата и услугата отказва да подпише маркер, ако тя липсва или е по-къса от 24 знака;
в изходния текст няма стойност по подразбиране.

Входните точки за удостоверяване са ограничени по честота до 30 обръщения за 15
минути от адрес. Ограничението е поставено само върху тях, а не глобално, защото
останалите маршрути са защитени с маркер и злоупотребата с тях изисква вече открадната
сесия.

\subsection{Валидиране на входа на границата}
\label{subsec:validation}

Всяко поле, което идва отвън, преминава през преобразуване, което го привежда към
известен вид и известен размер, преди да достигне до базата от данни
(Листинг~\ref{lst:validate}). Подходът е \emph{изграждане на разрешеното}, а не
премахване на забраненото: за състоянието на калкулатора се изброяват допустимите
имена на полета и се пренасят само те, така че произволно поле, изпратено от клиент,
не може да попадне в записа.

\begin{lstlisting}[caption={Изграждане на разрешеното вместо премахване на забраненото},label={lst:validate}]
function cleanInput(v) {
  if (!v || typeof v !== "object") return {};
  const allow = ["units", "type", "height", "levels", "span",
                 "overhang", "force", "factored", "capacity"];
  const out = {};
  for (const k of allow) if (v[k] !== undefined && v[k] !== null) out[k] = v[k];
  return out;
}
\end{lstlisting}

По същия начин етикетите се съкращават до 40 знака и до шест на брой с премахване на
повторенията, потребителските свойства до 60 знака за ключа, 500 за стойността и 50
двойки, а името на проекта до 120 знака. Границите не са произволни: те са избрани
така, че да не ограничават реалната употреба, но да поставят горна граница на
размера на един запис.

\subsection{Реализация на уеб клиента}
\label{subsec:webimpl}

Уеб клиентът е съвкупност от независими файла на стандартен JavaScript, всеки от
които се самоизпълнява в собствена област и излага в глобалната област най-много един
обект. Главният от тях, \code{app.js} с 1380 реда, съдържа целия калкулатор.

Състоянието е един обект \code{S} с деветте входни величини плюс избраните
допълнителни условия. Успоредно с него се води втори обект \code{selectedInputs} с
логически стойности, който отбелязва \emph{кои от величините потребителят е избрал
съзнателно}. Двата обекта не могат да бъдат слети: \code{S} винаги съдържа валидна
комбинация, защото се привежда след всяка промяна, докато \code{selectedInputs}
отговаря на въпроса дали тази валидна комбинация е поискана от човек.

Наборът от допустими стойности се извлича от самите данни при всяко преизчисляване, а
не се записва като отделна таблица (Листинг~\ref{lst:options}). Решението премахва
цял клас грешки: невъзможно е списъкът с допустими стойности да се разсинхронизира с
данните, защото той \emph{е} данните.

\begin{lstlisting}[caption={Допустимите стойности се извличат от самите данни},label={lst:options}]
function optionsFor() {
  if (S.type === "boulder") {
    return { heights: boulderHeights,
             spans: uniq(DATA.boulder.rows.map(r => r.a)).sort(num),
             overhangs: uniq(DATA.boulder.rows.map(r => r.x)).sort(num) };
  }
  var w = DATA.walls[wallKey(S.height, S.levels)];
  var spans = w ? uniq(w.rows.map(r => r.a)).sort(num) : [];
  var over  = w ? uniq(w.rows.map(r => r.x)).sort(num) : [];
  return { heights: wallHeights, spans: spans, overhangs: over };
}
\end{lstlisting}

Привеждането, чиято блок-схема е дадена на \figref{fig:clamp}, е реализирано в
\code{clampAndRender} (Листинг~\ref{lst:clamp}). Забележима е двойката действия при
всяка поправка: стойността се заменя с предпочитана по подразбиране, ако тя е налична,
и едновременно се отбелязва като неизбрана.

\begin{lstlisting}[caption={Привеждане на избора към допустима комбинация},label={lst:clamp}]
function clampAndRender() {
  var opt = optionsFor();
  if (opt.spans.indexOf(S.span) < 0) {
    S.span = opt.spans.includes(6) ? 6 : opt.spans[0];
    selectedInputs.span = false;
  }
  if (opt.overhangs.indexOf(S.overhang) < 0) {
    S.overhang = opt.overhangs.includes(1) ? 1 : opt.overhangs[0];
    selectedInputs.overhang = false;
  }
  renderControls();
  renderResults();
}
\end{lstlisting}

Текущият избор се запазва в локалното хранилище на браузъра при всяка промяна, така
че затварянето на страницата да не изисква повторно въвеждане. Записът съдържа само
входните величини, не и резултата, защото резултатът е производен.

\subsection{Изчислителното ядро}
\label{subsec:coreimpl}

Ядрото се свежда до три действия. \term{Намиране на реда} става чрез съставяне на ключ
от височината и броя нива и филтриране на редовете по разстояние между колоните и
наклон. \term{Избор на набор} чете полето \code{eu} или \code{us} според мерната
система, без преобразуване, по причината от §\ref{subsec:tables}. \term{Прилагане на
частните коефициенти} умножава стойността с коефициента за постоянен или за полезен
товар, но само когато потребителят е поискал изчислителни стойности.

Форматирането на числата зависи от мерната система: две значещи цифри след
десетичния знак за метрична и една за имперска. Разликата не е козметична, а следва
от порядъка на величините, килонютони срещу фунтове, при който еднакъв брой цифри би
означавал различна точност.

\subsection{Графичната схема}
\label{subsec:schematicimpl}

Схемата се съставя от модул за геометрия, който получава разстоянието между колоните,
наклона, височината и височините на нивата, и връща координатите на всички елементи:
контурите на катерачната повърхност, колоните на сградата, гредите, точките на
закрепване, стрелките на товарите, размерните линии и надписите. Модулът не докосва
документа; съставянето на елементите е задача на слоя на представянето.

Схемата поддържа увеличение и преместване с показалеца и с докосване. При тесен екран
това не е удобство, а необходимост: чертеж с размерни линии, смален до ширината на
телефон, става нечетлив, докато същият чертеж с превъртане остава използваем.

\subsection{Конфигуратор на закрепването}
\label{subsec:attachimpl}

Изборът на типов детайл за закрепване следва същия принцип на привеждане и явно
потвърждение като калкулатора (§\ref{subsec:core}), приложен към друго състояние.
Потребителят избира вида на подовата плоча или на носещата колона, а от списъка с
типови детайли, филтриран по този избор, посочва точно \emph{един}. Докато конкретен
детайл не е избран (\code{state.detail} е празно), точките на закрепване в схемата не
се изчертават като интерактивни и не носят стрелки на товари — вижда се само
очертанието на съоръжението. След избора точките на закрепване стават фокусируеми
елементи с роля на бутон (NFR3) и при посочване показват преглед на съответния типов
детайл. Избраният детайл се обозначава с промяна на надписа на бутона от
\emph{Select} на \emph{Selected} и с клас за визуално открояване, а не само с цвят,
за да остане различим и без възможност за разчитане на цвета. Изборът се пази в
локалното хранилище на браузъра, отделно за всяка от двете схеми на закрепване, и се
възстановява при следващо отваряне на страницата.

\subsection{Проверка на реализацията и получени резултати}
\label{subsec:method}

\subsubsection{Начин на проверка}
\label{subsec:howcheck}

Реализацията има три твърдения, които подлежат на проверка, и за всяко от тях
съществува различен по вид еталон.

\textbf{Числата са верни (NFR1).} Тук е налице нещо, което рядко съществува при
разработката на приложен софтуер: \term{оракул}, тоест независим източник, за който е
прието, че дава верния отговор~\cite{barr2015oracle}. Входното пространство е крайно и
малко, което позволява проверката да бъде изчерпателна, а не извадкова: 910 комбинации
за катерачни стени, 70 за боулдер стени и до дванадесет числа в метрична и дванадесет в
имперска мярка за всяка, тоест общо 26\,600 стойности. За всяка комбинация ядрото се
извиква през същия път, по който го извиква интерфейсът, а не чрез пряк достъп до
записа. Сравнението не преминава през преобразуване на единици по причината от
§\ref{subsec:tables}, а за комбинация без ред в еталона очакваният резултат е изричното
съобщение за липса: проверката се проваля и когато системата произведе число там,
където еталонът мълчи.

\textbf{Интерфейсът остава използваем при всяка ширина (NFR4).} За „добър интерфейс“
еталон няма, затова изискването е формулирано не като качество, а като механично
проверимо \emph{свойство}: нито един елемент не разширява страницата отвъд ширината на
изгледа. Страницата се зарежда в рамка (елемент \code{iframe}) с точно зададена ширина,
която образува собствен изглед, така че правилата за пречупване се изчисляват спрямо
нея. Обхожда се всеки елемент и се проверява дали дясната му граница надхвърля ширината
на изгледа. Изключват се елементите с неподвижно разположение, скритите елементи и
онези, чийто предшественик е съд с превъртане или отрязване, защото за тях излизането
извън границата е проектното намерение (§\ref{subsec:uidesign}). Отделно се сравнява
ширината на превъртане на документа с ширината на изгледа: без тази втора проверка
препълването би било скрито от измерването точно както е скрито от потребителя, тоест
дефектът щеше да изглежда поправен, защото е станал невидим. Проверяват се трите
състояния на калкулатора, които произвеждат различна разметка, и двете състояния на
удостоверяване; състоянията се достигат чрез програмно предизвикани събития върху
същите бутони, които потребителят натиска.

\textbf{Договорът на приложния програмен интерфейс се спазва.} Услугата се стартира
върху свободен порт от примката, а базата от данни е заместена с равностойна реализация
в паметта, така че изпитанието да не зависи от мрежа, от разрешаване на имена и от
външна услуга. Заявките се изпращат с истински обръщения по протокола, а бисквитките се
пренасят между тях, тоест изпитва се и транспортът на сесията.

\subsubsection{Съответствие с еталона}
\label{subsec:oracle}

Наборът с товарните таблици съществува в две копия, за услугата и за уеб клиента.
Структурното им сравнение дава съвпадение по всички 26\,600 стойности. Проверката е
евтина и е добавена, защото разминаването между тях би довело до най-неприятния възможен
дефект: два клиента, които показват различни числа за един и същ вход, без нито един от
тях да съобщава грешка.

Решението да се \emph{избира}, а не да се преобразува набор (§\ref{subsec:tables}) беше
прието по съображение, а след това проверено чрез измерване \tabref{tab:units}.

\begin{table}[htbp]
\centering
\caption{Възпроизводимост на имперския набор от метричния (13\,300 двойки)}
\label{tab:units}
{\fontsize{10}{12}\selectfont
\begin{tabular}{@{}L{8.4cm}rr@{}}
\toprule
\textbf{Правило за получаване} & \textbf{Съвпадения} & \textbf{Дял} \\
\midrule
точно преобразуване, $1$ kN $= 224{,}808943$ lbf & 0 & 0{,}00 \% \\
преобразуване със закръглен множител 225, закръглено до 0{,}1 lb & 7\,007 & 52{,}68 \% \\
точно преобразуване, закръглено до 0{,}1 lb & 95 & 0{,}71 \% \\
\addlinespace
в границите на 1 \% от точното преобразуване & 13\,260 & 100{,}00 \% \\
\bottomrule
\end{tabular}}
\end{table}

Наблюдаваното отношение между двата набора се движи между 224{,}570 и 225{,}359 при
медиана точно 225{,}000, тоест двата набора описват едни и същи физически величини и
разликите са от порядъка на закръглянето. Въпреки това \emph{нито едно} правило за
преобразуване не възпроизвежда имперския набор изцяло: най-доброто дава малко над
половината от стойностите, а най-голямото отклонение от него е 0{,}162 lb. Обяснението
е, че в изходната работна книга имперската колона е получена от незакръглени междинни
стойности, а не от показаните метрични. Изводът е практически: ако системата
преобразуваше единици вместо да избира набор, тя щеше да произведе числа, различни от
еталона при около 47 \% от входните комбинации, при това с малка разлика, тоест по
начин, който трудно би бил забелязан при оглед.

\subsubsection{Проверка на разположението}
\label{subsec:viewportres}

Проверката е проведена два пъти върху една и съща кодова база: веднъж върху изходното
състояние и веднъж след отстраняването на намерените дефекти. Изходното състояние дава
17 нарушения от 80 съчетания страница и ширина \tabref{tab:before}.

\begin{table}[htbp]
\centering
\caption{Хоризонтално препълване преди отстраняването на дефектите, в пиксели над ширината на изгледа}
\label{tab:before}
{\fontsize{10}{12}\selectfont
\begin{tabular}{@{}L{3.9cm}rrrrrrr@{}}
\toprule
\textbf{Страница} & \textbf{320} & \textbf{375} & \textbf{414} & \textbf{600} & \textbf{768} & \textbf{900} & \textbf{1024} \\
\midrule
Техническа документация & 681 & 626 & 587 & 401 & 233 & 101 & --- \\
Детайли за закрепване & 758 & 703 & 664 & 478 & 310 & 178 & 54 \\
Поверителност и бисквитки & 332 & 277 & 238 & 74 & --- & --- & --- \\
\bottomrule
\end{tabular}}
\end{table}

Легенда към \tabref{tab:before}: числото е разликата между ширината на превъртане на
документа и ширината на изгледа; „---“ означава липса на нарушение. Останалите пет
страници нямат нарушения при нито една проверена ширина. Съществено за оценката на
дефекта е, че той \emph{не се вижда}: стиловете задават отрязване по хоризонталата на
най-външния елемент, поради което страницата не се измества настрани, а част от
съдържанието просто липсва, без указание, че липсва. При страницата с детайлите за
закрепване при ширина 320 px това означава, че 70 \% от полезната ширина на съдържанието
са недостижими.

След отстраняването на дефектите проверката е повторена върху разширен набор от 38
ширини от 320 до 2560 px, включително двете страни на всяка точка на пречупване, за осем
страници, в три състояния на калкулатора и в две състояния на удостоверяване. Нарушения
не са установени.

\subsubsection{Систематизация на установените дефекти}
\label{subsec:taxonomy}

Седемнадесетте нарушения се свеждат до три първопричини, а огледът на екранните форми
добави още три вида дефект, които не пораждат препълване, но нарушават използваемостта
\tabref{tab:taxonomy}.

\begin{table}[htbp]
\centering
\caption{Класове установени дефекти}
\label{tab:taxonomy}
{\fontsize{10}{12}\selectfont
\begin{tabular}{@{}cL{5.0cm}L{3.1cm}L{2.6cm}@{}}
\toprule
\textbf{Клас} & \textbf{Първопричина} & \textbf{Проява} & \textbf{Открит чрез} \\
\midrule
Д1 & централно подравнен основен блок във флекс-контейнер приема вътрешната си минимална ширина & препълване на трите страници & обхождане \\
Д2 & пътечка \code{1fr} вместо \code{minmax(0,1fr)} в решетка & препълване, вложено в Д1 & обхождане \\
Д3 & пътечка с фиксиран минимум, по-широк от съда & препълване при тесен екран & обхождане \\
Д4 & неявно подреждане в решетка поставя елемент на собствен ред & загубена височина в заглавната лента & оглед \\
Д5 & прилепен елемент с отместване, изчислено за друга точка на пречупване & празнина под горния ръб & оглед \\
Д6 & правило за носител с невъзможно условие & мъртви стилове & преглед на изходния текст \\
\bottomrule
\end{tabular}}
\end{table}

\textbf{Д1} е първопричината на всичките 17 нарушения и е неочевидна. Тялото на
документа е флекс-контейнер с посока по колона, а основният блок на всяка страница е
централно подравнен чрез автоматични външни отстъпи. По спецификацията автоматичен
отстъп по напречната ос отменя разтягането на елемента до ширината на
контейнера~\cite{css_flexbox1}; вместо това елементът получава размер по съдържание,
тоест поне вътрешната си минимална ширина~\cite{css_sizing3}. Достатъчен е един широк
наследник, например таблица с изрична минимална ширина 620 px или чертеж в мащаб 1:1.
Отстраняването е едно правило, което задава изрична ширина 100 \% на основния блок, след
което наследниците с превъртане поемат излишъка (Приложение Г).

\textbf{Д2} се проявява едва след отстраняването на Д1 и илюстрира защо поправките
трябва да се проверяват повторно, а не да се приемат за успешни. Записът \code{1fr} е
съкращение на \code{minmax(auto, 1fr)}, при което автоматичният минимум е вътрешната
минимална ширина на съдържанието~\cite{mdn_grid}. Докато блокът-родител е бил по-широк
от изгледа, това не е било забележимо; след ограничаването му решетката е станала новото
място, през което съдържанието излиза. Останалите четири класа са съответно превантивно
открита пътечка с фиксиран минимум 290 px при екран 320 px (Д3), навигация през цялата
ширина, която изтласква бутоните за профил на собствен ред (Д4), прилепнала лента с
отместване, изчислено за друга точка на пречупване (Д5), и правило за носител с условие
\code{(min-width: 2000px) and (max-width: 1000px)}, което не може да бъде изпълнено и
чиито 87 реда стилове са премахнати (Д6).

Общата констатация е, че половината класове са откриваеми механично, а другата половина
изисква оглед. Първите три произтичат от правила за определяне на размер, при които
подразбирането е „по съдържание“, а намерението е „по съда“. От това следва практическо
заключение: автоматичната проверка на разположението не заменя огледа, а освобождава
време за него, като поема класа дефекти, който е най-многоброен и най-скучен за
откриване на ръка.

\subsubsection{Приложен програмен интерфейс и заплахи за валидността}
\label{subsec:apires}

Изпитанието на приложния програмен интерфейс преминава с 36 успешни проверки от 36 при
продължителност под една секунда (Приложение Д). Обхванати са състоянието на услугата,
отхвърлянето на неправилен адрес и на къса парола, създаването на профил с поставяне на
бисквитка, извличането на текущия профил със и без сесия, отхвърлянето на повторна
регистрация, входът с грешна и с вярна парола, издаването и приемането на маркер през
заглавие \code{Authorization}, пълният цикъл на проекта заедно със съхранените избори
на конфигуратора на закрепването, филтрирането и подреждането на списъка, обхватът на
достъпа до чужд проект и до чуждо запитване, и изходът от системата.

Резултатите имат четири известни ограничения. Проверката на интерфейса измерва едно
свойство, не качество: липсата на хоризонтално препълване не гарантира четливост,
достатъчен размер на целите за докосване или смислена подредба, и тези страни са
проверени чрез оглед, тоест описателно. Изпитанието на приложния програмен интерфейс
използва заместител на базата от данни, поради което зависещото от нея поведение се
възпроизвежда по описание. Еталонът е една версия на работната книга и резултатът не се
пренася автоматично към следваща редакция. Проверката е върху един двигател за
изобразяване, а различията между браузърите при изчисляване на вътрешната минимална
ширина са известен източник на разминавания и не са обхванати.

\subsection{Известни ограничения на реализацията}
\label{subsec:limitations}

\begin{enumerate}[topsep=2pt,itemsep=2pt,leftmargin=*]
\item \textbf{Изчислителният слой и слоят на представянето не са разделени по файлове}
      в приетия уеб клиент, а само по функции (§\ref{subsec:layers}). Това е цената на
      отказа от стъпка на компилация и означава, че ядрото на уеб клиента не може да
      бъде изпитано изолирано от слоя на представянето.
\item \textbf{Обобщението на проекта се съхранява отделно от входа}
      (§\ref{subsec:datamodel}). При промяна в товарните таблици обобщенията в списъка
      остават стари, докато проектът не бъде отворен и записан отново.
\item \textbf{Наборът с товарните таблици съществува в две копия} (за услугата и за
      уеб клиента), които не се съгласуват автоматично. Разминаване между тях би било
      открито от изпитанието по §\ref{subsec:oracle}, но не е предотвратено по
      конструкция.
\item \textbf{Документът \code{server/API.md} е остарял} в две точки: описва път
      \code{/api/projects/:id/questions}, докато реализираният е
      \code{/api/projects/:id/support-inquiries}, и описва извличане на списъка със
      запитвания, каквото не е реализирано.
\item \textbf{Един файл на уеб клиента не се зарежда от нито една страница}
      (\code{cookie-consent.js}), тоест лентата за съгласие за бисквитки е написана,
      но не е включена. Функционалността ѝ не е обхваната от изискванията и файлът е
      отбелязан тук, за да не бъде приет за действащ.
\end{enumerate}
